% Corrections relative to the earlier release (v0.5.0). This is the ONLY place values of that
% release appear, each marked superseded and mapped to the defects that supersede it. The defects
% are recorded in the project's publication-rigour review register and deferred-findings list,
% released with the code; the run-record amendments are in paper/data/run_record_amendments.json.

\section{Corrections relative to v0.5.0}\label{app:corrections}

This appendix concerns v0.5.0, the release of this work that preceded the pre-registered
study and is archived under the same record as this paper. All values reported in that earlier
release are superseded; none is carried into this version. Table~\ref{tab:corrections} maps
each superseded value to the defects behind it. The figures of the earlier release quoted in this appendix are read from its archived
record (added at commit \texttt{\resVFiveRecordCommit} and unchanged since), never typed. The
narrative figures (the Norman split, the share of unannotated cells in one task, the almost
disjoint task sets, the four of five roles and the four pooled constructs) are quoted from the
internal audit of the earlier release that found them, and each quote is checked against the
audit's record. The
files released with the earlier release do not describe a single experiment, and several of its
values were computed on inputs that a code defect had altered. The defects, whose
audit findings are released with the code, were the following.

\begin{description}[leftmargin=1.2em,itemsep=1pt,font=\normalfont\bfseries]
  \item[Disjoint task sets.] The trainer and lifecycle files covered almost disjoint task
    sets (2 shared of \resVFiveNTasks{} each), produced by two processes whose task draws
    diverged because the stratification was salted per process; statistics joining them
    were computed across that intersection.
  \item[Seeds not propagated.] The seed was recorded but not passed to the lifecycle trainer
    or the LLM cache key, so the three seeds per task were byte-identical copies.
  \item[Step source not recorded.] Lifecycle steps did not record whether a proposal came
    from an LLM or the rule-based fallback; four of five roles emitted their schema default
    throughout.
  \item[Doublet labels misparsed.] Norman doublet labels were tested with the wrong delimiter,
    so the doublet stratum was empty and each doublet task was trained and scored with a
    randomly chosen gene as its target.
  \item[Datasets not pinned.] No dataset digest was pinned; the fetchers were described as
    verifying digests but trusted any cached file.
  \item[Unresolved Norman labels.] Norman labels that did not resolve to a symbol were given a
    random target and entered the training set of every Norman task not holding them out.
  \item[Pooled Adamson constructs.] The Adamson task \textit{ATF6} pooled four distinct
    constructs (a single, two doubles and a triple) into one population.
  \item[Missing annotations decoded.] Adamson cells with a missing annotation were decoded as
    the file's last category; the \textit{YIPF5} population was about 91\% unannotated cells,
    and both \textit{YIPF5} and \textit{ZNF326} entered training for other Adamson tasks and
    were lifecycle held-out tasks.
  \item[Gene selection saw held-out cells.] Highly variable genes were selected on all cells,
    including the held-out perturbation's cells.
  \item[Reporting defects.] Criteria that could not be evaluated were dropped from the tally;
    the language models that served the agents were misstated and not recorded; the
    documented Norman design contradicted the released files; and the provenance record held no code revision, seed or
    dataset digest.
\end{description}

\begin{table}[h]
\small
\centering
\begin{tabular}{p{0.44\linewidth}p{0.24\linewidth}p{0.22\linewidth}}
\toprule
Superseded value or claim of the earlier release & Superseded because of & Replacement \\
\midrule
Adamson median of the lowest mean squared deviation (MSD) over configurations, $\resVFiveAdamsonMedian$, reported as meeting its criterion
  & disjoint task sets, pooled Adamson constructs, missing annotations decoded, gene selection saw held-out cells
  & Adamson attainable-error hypothesis, \resHOneVerdict \\
Norman median of the lowest MSD over configurations, $\resVFiveNormanMedian$, reported as meeting its criterion
  & disjoint task sets, doublet labels misparsed, unresolved Norman labels, gene selection saw held-out cells
  & Norman attainable-error hypothesis, \resHTwoVerdict \\
Norman $n=15$ described as $15$ singletons $+\,5$ doublets (the released files hold $7+8$)
  & reporting defects, doublet labels misparsed
  & $\resHTwoSingletonN+\resHTwoDoubletN$ asserted at preflight \\
$\resVFiveNTrainerRuns$ trainer records ($\resVFiveNTasks$ tasks $\times$ $\resVFiveConfigsPerTask$ configurations)
  & disjoint task sets, doublet labels misparsed
  & $\resNTasks$ tasks $\times$ \resNTrainerRecordsPerTask{} records (\resNDistinctConfigs{} distinct configurations) \\
$\resVFiveNLifecycleRuns$ lifecycle runs over $\resVFiveNTasks$ held-out tasks
  & disjoint task sets, seeds not propagated
  & $\resNTasks$ tasks $\times$ $3$ seeds, same task set as the trainer \\
Architect backbone entropy $\resVFiveBackboneEntropy$ nats; the transformer in $\resVFiveScgptSharePct\%$ of
  $\resVFiveNPicks$ picks; entropy of the count of highly variable genes (HVGs) $\resVFiveHVGEntropy$ nats
  & disjoint task sets, seeds not propagated, step source not recorded, reporting defects, missing annotations decoded
  & backbone-choice hypothesis, \resHThreeVerdict; HVG-count entropy \resHVGEntropy{} nats \\
Rank correlations $\rho$ of the Task Difficulty Index (TDI) and its components with error, and the transfer $\rho$, reported as not available
  & reporting defects
  & trace-feature ranking and cross-dataset transfer hypotheses \\
``Two of three pre-registered gates pass''
  & reporting defects & tally out of five \\
``\textsc{Nemotron-30B} throughout''
  & reporting defects & pinned Anthropic roster, served model asserted per call \\
Spend \$\resVFiveSpend{} over \resVFiveGPUHours{} GPU-hours; \$\resVFiveBudgetCap{} quoted as the envelope
  & disjoint task sets, reporting defects & LLM \$\resSpendLLM, GPU \$\resSpendGPU{} (\resGPUHours{} A100-hours), total \$\resSpendTotal \\
Fetchers described as verifying digests & datasets not pinned & four digests pinned, failing closed \\
Byte-equivalent re-run from ``the same code revision'' & reporting defects & re-run at the
  recorded revision and configuration \\
\bottomrule
\end{tabular}
\caption{Each superseded value or claim of the earlier release (v0.5.0), the defects that
supersede it, named by the headings of the list that precedes this table, and its replacement
in this version. The hypotheses named in the last column are those of
Table~\ref{tab:gates}.}
\label{tab:corrections}
\end{table}

\subsection*{Amendments to the run record}\label{app:rr}
\begin{sloppypar}
The run record was amended twice after the run; both amendments are logged with the code, and
the hashes of every file before and after are released with the data. The first amendment
(\resRROneDate, decided by the author) removed the record of where the LLM API
key was stored (the name of the secret store, an organisation project, an environment, a home
project and a cross-project flag; no secret value) from \resRROneFileCount{} run files, the run
record and the three run manifests, because it supports no claim. The SHA-256 prefix of the run
manifest, \resManifestShaSuperseded{} as first published, is superseded by \resManifestSha. No
result changes. The code at the run's revision records the field; the code was changed so that
it no longer does, and a re-run as now documented does not record it. The second amendment
(\resRRTwoDate) changed no file or hash. It corrected three notes of the first amendment, which
had wrongly said that the code still recorded the field and that this was out of scope, and left
the first amendment's own record unchanged.
\end{sloppypar}
